% ----------------------------------------------------------------------
%  Введение к главе 2
% ----------------------------------------------------------------------
В главе~1 кристалл был идеальным: атомы стоят строго в узлах решётки,
и картина бесконечно повторяется. Реальный кристалл устроен иначе~— в
нём всегда есть нарушения порядка, \term{дефекты}. Их немного: даже в
сильно деформированном металле в дефектных местах находится лишь малая
доля атомов. Но именно дефекты определяют многие важнейшие свойства
материала.

\begin{glavnoe}[Зачем изучать дефекты]
Свойства кристалла делятся на две группы.
\emph{Структурно-нечувствительные} (плотность, упругие модули, температура
плавления) почти не зависят от дефектов~— их можно рассчитать для
идеальной решётки. \emph{Структурно-чувствительные} (прочность и
пластичность, скорость диффузии, остаточное электросопротивление,
поведение под облучением) определяются именно дефектами. Поэтому
материаловед управляет свойствами, управляя дефектами.
\end{glavnoe}

\subsubsection*{Классификация дефектов по размерности}

Дефекты удобно различать по тому, во скольких измерениях они протяжённы
(рис.~\ref{fig:defect-zoo}, табл.~\ref{tab:defect-classes}).

\begin{table}[ht]
\centering
\caption{Геометрическая классификация дефектов кристаллического строения}
\label{tab:defect-classes}
\small
\begin{tabularx}{\textwidth}{@{}l>{\raggedright\arraybackslash}X>{\raggedright\arraybackslash}X@{}}
\toprule
\textbf{Класс} & \textbf{Размеры} & \textbf{Примеры} \\
\midrule
Точечные (0D) & порядка межатомного расстояния во всех трёх
  направлениях & вакансия, межузельный атом, примесный атом, их
  комплексы \\
Линейные (1D) & малы в двух направлениях, а в третьем сравнимы с
  размером кристалла & дислокации, дисклинации, цепочки вакансий \\
Поверхностные (2D) & малы только в одном направлении &
  поверхность кристалла, границы зёрен и фаз, двойниковые границы,
  дефекты упаковки, антифазные границы \\
Объёмные (3D) & велики во всех трёх направлениях & поры, трещины,
  включения и выделения второй фазы \\
\bottomrule
\end{tabularx}
\end{table}

\begin{figure}[ht]
\centering
\begin{subfigure}[b]{0.24\textwidth}\centering
\begin{tikzpicture}[scale=0.42]
  \foreach \x in {0,...,6}{\foreach \y in {0,...,6}{
    \pgfmathtruncatemacro{\code}{10*\x+\y}
    \ifnum\code=24 \node[vac=3.4mm] at (\x,\y) {};
    \else\ifnum\code=42 \node[atomB=4.4mm] at (\x,\y) {};
    \else \node[atom=3.4mm] at (\x,\y) {};\fi\fi}}
  \node[atomI=3.0mm] at (4.5,4.5) {};
  \node[atom=2.0mm, fill=atomLine!70] at (1.5,0.5) {};
\end{tikzpicture}
\caption{точечные}
\end{subfigure}\hfill
\begin{subfigure}[b]{0.24\textwidth}\centering
\begin{tikzpicture}[scale=0.42]
  \def\edgeAtomR{0.17}
  \begin{scope}\clip (-3.3,-3.3) rectangle (3.3,3.3);
  \EdgeDislLattice{4}{4}{4}\end{scope}
  \dislT{(0,-0.05)}{0}
\end{tikzpicture}
\caption{линейный}
\end{subfigure}\hfill
\begin{subfigure}[b]{0.24\textwidth}\centering
\begin{tikzpicture}[scale=0.42]
  \begin{scope}
    \clip (-3.3,-3.3) rectangle (0,3.3);
    \begin{scope}[rotate=14]
      \foreach \x in {-6,...,2}{\foreach \y in {-6,...,6}{
        \node[atom=3.2mm] at (\x,\y) {};}}
    \end{scope}
  \end{scope}
  \begin{scope}
    \clip (0,-3.3) rectangle (3.3,3.3);
    \begin{scope}[rotate=-14]
      \foreach \x in {-2,...,6}{\foreach \y in {-6,...,6}{
        \node[atom=3.2mm, fill=atomA!55] at (\x,\y) {};}}
    \end{scope}
  \end{scope}
  \draw[fmRed, dashed, line width=0.8pt] (0,-3.4) -- (0,3.4);
\end{tikzpicture}
\caption{поверхностный}
\end{subfigure}\hfill
\begin{subfigure}[b]{0.24\textwidth}\centering
\begin{tikzpicture}[scale=0.42]
  \foreach \x in {-3,...,3}{\foreach \y in {-3,...,3}{
    \pgfmathsetmacro{\rr}{(\x-0.3)^2+(\y-0.2)^2}
    \pgfmathparse{\rr<3.5 ? 1 : 0}
    \ifnum\pgfmathresult=0 \node[atom=3.4mm] at (\x,\y) {};\fi}}
  \draw[atomLine, dashed] (0.3,0.2) circle (1.95);
\end{tikzpicture}
\caption{объёмный}
\end{subfigure}
\caption{Дефекты разной размерности в двумерной модели кристалла:
\textit{а}~— вакансия (пунктир), межузельный атом (красный),
крупный атом замещения (оранжевый) и мелкий атом внедрения (тёмный);
\textit{б}~— краевая дислокация (оранжевым выделена «лишняя» атомная
полуплоскость, значок $\perp$~— линия дислокации);
\textit{в}~— граница между двумя зёрнами с разной ориентацией;
\textit{г}~— пора}
\label{fig:defect-zoo}
\end{figure}

Особое место занимают \term{дислокации}. Это не просто «ещё один»
линейный дефект: с их помощью объясняются пластическая деформация и
упрочнение металлов, рост кристаллов, многие фазовые превращения. Поэтому
дислокациям отведена бо́льшая часть главы.

\begin{figure}[ht]
\centering
\begin{tikzpicture}[
  box/.style={draw=#1, fill=#1!8, rounded corners=2pt, align=left,
    font=\small, inner sep=5pt, text width=60mm, line width=0.6pt},
  hdr/.style={font=\sffamily\bfseries\small, text=#1},
  flow/.style={-{Stealth[length=2mm]}, line width=0.6pt, fmGray}]
\node[box=fmBlue] (p) {
\textbf{2.1.1}~Вакансии и межузельные атомы\\
\textbf{2.1.2}~Энергия образования, концентрация, подвижность\\
\textbf{2.1.3}~Комплексы точечных дефектов\\
\textbf{2.1.4}~Источники, стоки, отжиг\\
\textbf{2.1.5}~Как получить много дефектов\\
\textbf{2.1.6}~Как дефекты измеряют};
\node[hdr=fmBlue, above=1pt of p.north west, anchor=south west] {2.1. Точечные дефекты};
\node[box=fmAmber, right=17mm of p] (d) {
\textbf{2.2.1}~Краевая и винтовая дислокации\\
\textbf{2.2.2}~Скольжение и переползание\\
\textbf{2.2.3}~Поле напряжений, энергия, силы\\
\textbf{2.2.4}~Откуда берутся дислокации\\
\textbf{2.2.5}~Скопления и стенки\\
\textbf{2.2.6--2.2.7}~Дислокации в ГЦК, ОЦК, ГПУ};
\node[hdr=fmAmber, above=1pt of d.north west, anchor=south west] {2.2. Дислокации};
\draw[flow, {Stealth[length=2mm]}-{Stealth[length=2mm]}] (p.east) -- node[above, font=\scriptsize, text=fmGray, align=center]{взаимо-\\действие} (d.west);
\end{tikzpicture}
\caption{Карта главы. Точечные дефекты и дислокации взаимодействуют:
вакансии позволяют дислокациям переползать, примеси тормозят дислокации,
а движение дислокаций само порождает точечные дефекты}
\label{fig:ch2-map}
\end{figure}
